% !TeX encoding = UTF-8
% !TeX program = xelatex
\documentclass[degree=bachelor,language=chinese,fontset=fandol]{thuthesis}

\usepackage{booktabs,longtable,array,xcolor,hyperref,enumitem}
\hypersetup{colorlinks=true,linkcolor=black,urlcolor=blue}
\setlist[itemize]{leftmargin=2em,itemsep=0.25em}
\newcommand{\yes}{\textcolor{green!50!black}{已生成}}
\newcommand{\no}{\textcolor{red!70!black}{未观察到}}
\newcommand{\na}{\textcolor{gray}{不适用}}

\begin{document}
\begin{titlepage}
  \centering
  \vspace*{2.3cm}
  {\zihao{1}\bfseries MLJAR 与 TabularX 完整工具包对比\par}
  \vspace{1.2cm}
  {\zihao{3}House Prices 数据集实验报告\par}
  \vspace{2.2cm}
  {\zihao{4}TabularX 课程项目\par}
  \vfill
  {\zihao{-4}实验日期：2026 年 9 月 29 日\par}
  \vspace{1.2cm}
\end{titlepage}

\frontmatter
\tableofcontents
\mainmatter
\chapter{实验概述}
\section{数据集与任务}
\textbf{数据集：House Prices（OpenML dataset 42165（house\_prices），回归）。}
House Prices 数据集包含 1,460 套住宅的 80 个原始描述特征，目标是预测住宅销售价格 SalePrice；标识列 Id 不进入模型。数据按固定随机种子 42 划分为 80\% 开发集和 20\% 外部测试集，外部测试集共 292 条记录。

\section{比较范围与原则}
实验比较 MLJAR 与 TabularX 两个完整工具包，范围包括预处理、模型搜索与选择、解释、公平性适用性、可信审计、模型持久化、报告和证据交付。双方使用同一数据划分，并运行 Random Forest、XGBoost、LightGBM 和 CatBoost。外部测试集不参与训练、调参和选模，只在最终模型锁定后评估一次。

性能指标只描述本次固定配置。主要比较对象是实际生成并能够核对的解释性与可信性工件。

\chapter{实验结果}
\section{外部测试集表现}
\begin{table}[htbp]
  \centering
  \small
  \caption{共同外部测试集上的模型表现}
  \begin{tabular}{llccc}
    \toprule
    工具 & 最终模型 & MAE & RMSE & $R^2$ \\
    \midrule
    TabularX & LightGBM & 17990.99 & 29635.38 & 0.8855 \\
    MLJAR & Ensemble & 16106.61 & 27228.68 & 0.9033 \\
    \bottomrule
  \end{tabular}
\end{table}

本次固定配置下，MLJAR 的 MAE（16106.61）和 RMSE（27228.68）低于 TabularX，且 $R^2$ 为 0.9033。该结果只描述本次小规模搜索，不构成普遍性能排名。

\section{可解释性与可信性工件盘点}
\begin{longtable}{>{\raggedright\arraybackslash}p{0.27\textwidth}>{\centering\arraybackslash}p{0.14\textwidth}>{\centering\arraybackslash}p{0.14\textwidth}>{\raggedright\arraybackslash}p{0.34\textwidth}}
  \caption{本次运行实际观察到的工件}\label{tab:artifacts}\\
  \toprule
  工件 & TabularX & MLJAR & 说明 \\
  \midrule
  \endfirsthead
  \toprule
  工件 & TabularX & MLJAR & 说明 \\
  \midrule
  \endhead
  最佳模型及可重载文件 & \yes & \yes & 两者均保存最终模型；TabularX 另行核验模型接收原始表格后的预测一致性。 \\
  候选模型与验证分数 & \yes & \yes & TabularX 保存 model\_scores.json；MLJAR 保存 leaderboard.csv 或参数记录。 \\
  预处理记录 & \yes & \yes & 两者均提供本次运行可查阅的预处理记录。 \\
  特征重要性 & \yes & \yes & 两者均生成特征重要性文件或图。 \\
  SHAP & \yes & \yes & 两者本次均观察到 SHAP 工件；覆盖范围以实际模型目录为准。 \\
  ALE & \no & \no & 本次固定配置下，两者均未观察到 ALE 工件。 \\
  学习曲线或调参轨迹 & \yes & \yes & TabularX 输出调参轨迹；MLJAR 输出学习曲线。 \\
  混淆矩阵 & \na & \na & 分类任务用于核对各类别预测错误；回归任务不适用。 \\
  公平性 & \na & \na & 数据集没有在本实验中声明敏感属性，因此公平性工件标记为不适用。 \\
  泄露、划分与评估协议审计 & \yes & \no & TabularX 将数据风险、划分完整性和测试集访问顺序写入审计证据。 \\
  漂移与群体覆盖 & \yes & \no & TabularX 检查开发集与外部测试集的分布差异及已声明群体覆盖。 \\
  重载一致性 & \yes & \no & TabularX 保存 reload\_check.json，核对持久化前后预测是否一致。 \\
  交付状态与失败原因 & \yes & \no & TabularX 保存 delivery.json，区分已生成、失败、缺失、未启用与不适用。 \\
  原生 HTML 报告 & \yes & \no & TabularX 生成汇总 HTML；MLJAR Explain 本次主要输出逐模型 Markdown 与图片。 \\
  \bottomrule
\end{longtable}

“未观察到”只表示本次配置和显式调用范围内没有找到对应工件，不外推为软件在任何条件下都不支持。

\chapter{工具优缺点与结果解读}
\section{TabularX}
\textbf{优点：}将测试集隔离、泄露与漂移审计、模型重载核验、报告生成和交付状态检查放入同一次运行的证据链。

\textbf{不足：}本次未观察到 ALE；当前小规模配置只能说明流程可运行，不能代表最优模型性能。

\section{MLJAR}
\textbf{优点：}Explain 模式能够生成逐模型 Markdown、学习曲线、混淆矩阵、特征重要性与 SHAP 等常规解释工件。

\textbf{不足：}本次运行目录中未观察到泄露与划分审计、漂移与群体覆盖、重载一致性、交付状态以及原生汇总 HTML 工件。

\chapter{结论}
两种工具均完成主流模型训练和常规解释工件生成。MLJAR 的逐模型说明较成熟；TabularX 在常规模型报告之外补充可信审计、测试隔离、重载核验和交付状态证据。后续数据集继续复用相同的数据隔离、工件状态词和结果边界。

\backmatter
\chapter{来源}
\begin{itemize}
  \item 数据来源：OpenML dataset 42165（house\_prices）
  \item MLJAR supervised：\url{https://github.com/mljar/mljar-supervised}
  \item THUThesis：\url{https://github.com/tuna/thuthesis}
\end{itemize}
\end{document}
